\documentclass[11pt]{article}

\usepackage[a4paper,margin=1in]{geometry}
\usepackage{amsmath,amssymb}
\usepackage{hyperref}

\hypersetup{
    colorlinks=true,
    linkcolor=blue,
    citecolor=blue,
    urlcolor=blue
}

\title{Law of a Random Variable}
\author{}
\date{}

\begin{document}

\maketitle

\section*{Definition}

Let $(\Omega,\mathcal F,\mathbb P)$ be a probability space and let
\[
X:\Omega\to E
\]
be a random variable taking values in a measurable space $(E,\mathcal E)$.
The law, or distribution, of $X$ is the probability measure on $E$ defined by
\[
\mathcal L(X)(A)
=
\mathbb P(X\in A)
=
\mathbb P(X^{-1}(A)),
\qquad A\in\mathcal E.
\]
Equivalently,
\[
\mathcal L(X)=X_\#\mathbb P.
\]

\section*{Time-Dependent Random Variables}

If $(X_t)_{t\ge 0}$ is a time-dependent random variable or stochastic process,
then the law of $X_t$ is
\[
\mathcal L(X_t)=(X_t)_\#\mathbb P.
\]
It describes the probability distribution of the random state at time $t$.

\end{document}
